\documentclass[11pt]{article}
\usepackage[a4paper,margin=24mm]{geometry}
\usepackage{amsmath,amssymb,amsthm,mathtools}
\usepackage{graphicx}
\usepackage{microtype}
\usepackage{xcolor}
\usepackage{enumitem}
\usepackage[hidelinks]{hyperref}

\newtheorem{theorem}{Theorem}
\newtheorem{lemma}[theorem]{Lemma}
\newtheorem{remark}[theorem]{Remark}
\newcommand{\E}{\mathbb E}
\newcommand{\Q}{\mathbb Q}
\newcommand{\cS}{\mathcal S}
\newcommand{\norm}[1]{\left\lVert #1\right\rVert}
\newcommand{\Tstar}{\mathsf T}

\title{A nonseparable strong-random extension of the\\
Gregu\v{s}--\'{C}iri\'{c} and Hardy--Rogers fixed-point theorems}
\author{Candidate solution packet for arXiv:1706.01634}
\date{11 August 2026}

\begin{document}
\maketitle

\begin{center}
\fbox{\parbox{0.92\textwidth}{\textbf{Status.}
Candidate full affirmative answer, likely valid, under the canonical
nonseparable formulation using strong (Bochner) random variables and strong
random sections.  Human verification is requested.}}
\end{center}

\begin{abstract}
Saipara, Kumam, and Cho ask whether two random fixed-point theorems in their
paper extend from separable to nonseparable Banach spaces.  We give an
affirmative extension to every real Banach space after making the standard
nonseparable measurability choice: random variables and constant operator
sections are strongly measurable.  The proof closes a countable family of
strong random variables under the operator and rational linear combinations.
Fibrewise, its closure is a separable invariant Banach subspace, so the
deterministic fixed-point theorem applies.  Countably selected approximate
residual minimizers converge to the fixed point by explicit residual error
bounds, proving its strong measurability.  The argument handles both the
Gregu\v{s}--\'{C}iri\'{c} condition and the five-term Hardy--Rogers condition.
\end{abstract}

\section{Source question}

The source is P.~Saipara, P.~Kumam, and Y.~J.~Cho,
\emph{Random fixed point theorems for Hardy--Rogers self-random operators
with applications to random integral equations}, arXiv:1706.01634v1 (2017),
published in \emph{Stochastics} 90 (2018), 297--311 \cite{source}.
On printed page 12, after the application theorem, it asks:

\begin{center}
\includegraphics[width=0.98\textwidth]{figures/open_problem_crop.png}
\end{center}

\begin{quote}
\textbf{Open Problem.} Can Theorems 1.3 and 3.2 be generalized to
non-separable Banach spaces?
\end{quote}

Theorem 1.3 is the random Gregu\v{s}--\'{C}iri\'{c} theorem quoted from
Saha--Ganguly \cite{saha}; Theorem 3.2 is the source paper's random
Hardy--Rogers theorem.  The source uses separability to identify its several
measurability notions.  In a nonseparable space those notions separate, so a
generalization must say which one is intended.

\section{Definitions and the precise extension}

Let $(\Omega,\mathcal F,\mu)$ be a complete probability space and let $X$ be
an arbitrary real Banach space.  An $X$-valued map is \emph{strongly
measurable} if it is the almost-sure norm limit of finite-valued simple maps.
Write $\cS(X)$ for the strong random variables modulo almost-sure equality.

A map $T:\Omega\times X\to X$ is a \emph{strong continuous random operator}
if
\begin{enumerate}[label=(S\arabic*),leftmargin=2.6em]
\item $\omega\mapsto T(\omega,x)$ is strongly measurable for every fixed
      $x\in X$;
\item $x\mapsto T(\omega,x)$ is norm-continuous for almost every $\omega$.
\end{enumerate}
This definition agrees with the source's continuous-random-operator
convention when $X$ is separable.

We use either of the following almost-sure conditions, imposed for every
$\xi,\eta\in\cS(X)$.

\paragraph{Gregu\v{s}--\'{C}iri\'{c} condition (GC).}
There are measurable $a,b,c:\Omega\to[0,\infty)$ with
\[
 0<a<1,\qquad a+b=1,
 \qquad c\le \frac{4-a}{8-a}
 \quad\text{a.s.},
\]
such that, abbreviating $T_\omega x=T(\omega,x)$,
\begin{align}
 \norm{T_\omega\xi-T_\omega\eta}
 &\le a\max\!\left\{\norm{\xi-\eta},
 c\bigl(\norm{\xi-T_\omega\eta}+\norm{\eta-T_\omega\xi}\bigr)\right\}
 \notag\\[-2mm]
 &\hspace{13mm}+
 b\max\!\left\{\norm{\xi-T_\omega\xi},
                  \norm{\eta-T_\omega\eta}\right\}.
 \tag{GC}
\end{align}
All values in (GC) are evaluated at the same $\omega$.

\paragraph{Hardy--Rogers condition (HR).}
There are measurable $\alpha_i:\Omega\to[0,\infty)$ such that
$\sum_{i=1}^{5}\alpha_i<1$ a.s. and
\begin{align}
 \norm{T_\omega\xi-T_\omega\eta}
 &\le \alpha_1\norm{\xi-\eta}
 +\alpha_2\norm{\xi-T_\omega\xi}
 +\alpha_3\norm{\eta-T_\omega\eta} \notag\\
 &\quad+\alpha_4\norm{\xi-T_\omega\eta}
 +\alpha_5\norm{\eta-T_\omega\xi}.
 \tag{HR}
\end{align}

\begin{theorem}[Nonseparable strong-random extension]\label{thm:main}
Let $X$ be any real Banach space and let $T:\Omega\times X\to X$ be a
strong continuous random operator.  If either \emph{(GC)} or \emph{(HR)}
holds for every pair in $\cS(X)$, then there is a unique
$u\in\cS(X)$ such that
\[
 T(\omega,u(\omega))=u(\omega)\qquad\text{for almost every }\omega.
\]
Consequently, both theorems named in the source question extend to arbitrary,
possibly nonseparable, Banach spaces in the strong-random formulation.
\end{theorem}

\section{Proof intuition}

The obstacle is not deterministic fixed-point geometry.  It is the loss of a
single countable dense subset with which to synchronize null sets and prove
measurability.  We create the needed countability from the operator itself.
Start at zero and repeatedly add all rational linear combinations and all
operator images.  This produces countably many strong random variables
$d_1,d_2,\ldots$.  At each outcome, their closed span is separable and
invariant under that fibre of $T$.  A countable intersection synchronizes the
contractive inequalities on this random subspace, where the deterministic
theorem supplies a fixed point.

The fixed point is not selected abstractly.  Among the $d_k$, choose the first
one whose residual $\norm{d_k-Td_k}$ is below $1/n$.  This is a measurable
countable choice.  Explicit estimates show that small residual forces small
distance to the fixed point, so these choices converge strongly to it.

This detour is essential for treating both source theorems uniformly.  Raw
Picard iteration need not converge in the Gregu\v{s}--\'{C}iri\'{c} case:
$T=-I$ satisfies (GC), has the unique fixed point $0$, and its nonzero Picard
orbits oscillate.

\section{Proof}

\subsection{Strong composition}

\begin{lemma}[Composition lemma]\label{lem:composition}
If $T$ is a strong continuous random operator and
$\xi\in\cS(X)$, then
\[
 (\Tstar\xi)(\omega):=T(\omega,\xi(\omega))
\]
belongs to $\cS(X)$.
\end{lemma}

\begin{proof}
Choose finite-valued simple $\xi_m$ with $\xi_m\to\xi$ almost surely in
norm.  On each member of the finite measurable partition defining $\xi_m$,
$\Tstar\xi_m$ equals a constant section $T(\cdot,x)$; hence
$\Tstar\xi_m$ is strongly measurable.  Fibre continuity gives
$\Tstar\xi_m\to\Tstar\xi$ almost surely.  Strong measurability is closed under
almost-sure norm limits.
\end{proof}

\subsection{A countable random invariant hull}

Set $\mathcal D_0=\{0\}$.  Recursively let
\[
 \mathcal D_{n+1}
 =\operatorname{span}_{\Q}
   \bigl(\mathcal D_n\cup\{\Tstar\xi:\xi\in\mathcal D_n\}\bigr),
 \qquad
 \mathcal D=\bigcup_{n\ge0}\mathcal D_n.
 \tag{1}
\]
Each $\mathcal D_n$ is countable, and Lemma~\ref{lem:composition} shows that
every member is strongly measurable.  Enumerate $\mathcal D$ as
$(d_k)_{k\ge1}$.  For each $\omega$, define
\[
 Y_\omega=\overline{\{d_k(\omega):k\ge1\}}.
 \tag{2}
\]
Because the point set in (2) is rational-linear, $Y_\omega$ is a closed real
linear subspace of $X$, hence a separable Banach space.  Construction (1) and
fibre continuity imply
\[
 T_\omega(Y_\omega)\subseteq Y_\omega.
 \tag{3}
\]

For every pair $(j,k)$ the assumed random inequality holds almost surely for
$(d_j,d_k)$.  Intersect these countably many full-measure events with the
events on which the coefficient restrictions and fibre continuity hold.  On
the resulting full-measure set $\Omega_0$, the relevant inequality holds on
all pairs $d_j(\omega),d_k(\omega)$ and therefore, by continuity, on all pairs
in $Y_\omega$.

If (GC) holds, the deterministic Gregu\v{s}--\'{C}iri\'{c} theorem
\cite{ciric2000} applied to the closed convex set $Y_\omega$ gives a unique
fixed point $u(\omega)\in Y_\omega$.  If (HR) holds, put
$q(\omega)=\sum_i\alpha_i(\omega)<1$.  The right-hand side of (HR) is at most
$q(\omega)$ times the maximum of the five usual \v{C}iri\'{c} distances, so
$T_\omega|_{Y_\omega}$ is a quasicontraction.  The deterministic
quasicontraction theorem \cite{ciric1974} again gives a unique fixed point
$u(\omega)\in Y_\omega$.

At this stage no measurability of $u$ has been assumed or inferred.

\subsection{Residual stability}

\begin{lemma}[Residual error bounds]\label{lem:residual}
Fix $\omega\in\Omega_0$, suppress it from the notation, and let $u=Tu$ be
the fixed point in $Y_\omega$.  For $x\in Y_\omega$ set
\[
 r=\norm{x-u},\qquad \delta=\norm{x-Tx},\qquad d=\norm{Tx-u}.
\]
Under (GC),
\[
 r\le C_{GC}\delta,
 \quad
 C_{GC}:=\max\left\{\frac{1+b}{b},
 \frac{1-ac+b}{1-2ac}\right\}<\infty.
 \tag{4}
\]
Under (HR),
\[
 r\le C_{HR}\delta,
 \quad
 C_{HR}:=\frac{1-\alpha_5+\alpha_2}
 {1-\alpha_1-\alpha_4-\alpha_5}<\infty.
 \tag{5}
\]
\end{lemma}

\begin{proof}
In the (GC) case, substituting $y=u$ gives
\[
 d\le a\max\{r,c(r+d)\}+b\delta,
 \qquad r\le\delta+d.
 \tag{6}
\]
If the maximum in (6) is $r$, then
$r\le\delta+ar+b\delta$, so
$r\le(1+b)\delta/(1-a)=(1+b)\delta/b$.
If the maximum is $c(r+d)$, then
\[
 (1-ac)d\le ac,r+b\delta.
\]
Combining this with $r\le\delta+d$ yields
$(1-2ac)r\le(1-ac+b)\delta$.  The source restriction implies
$c\le(4-a)/(8-a)<1/2$, so both denominators in (4) are positive.

In the (HR) case, substitution of $u=Tu$ gives
\[
 (1-\alpha_5)d
 \le(\alpha_1+\alpha_4)r+\alpha_2\delta.
\]
Combining with $r\le\delta+d$ gives
\[
 (1-\alpha_1-\alpha_4-\alpha_5)r
 \le(1-\alpha_5+\alpha_2)\delta,
\]
which is (5).  Its denominator is positive because
$\sum_i\alpha_i<1$.
\end{proof}

\subsection{Measurable approximation and conclusion}

For $k\ge1$, define the measurable residual
\[
 \rho_k(\omega)=
 \norm{d_k(\omega)-T(\omega,d_k(\omega))}.
\]
It is measurable because both terms are strongly measurable.  For each
$n\ge1$, on the event that an index with $\rho_k<1/n$ exists, let
$k_n(\omega)$ be the least such index and set
\[
 z_n(\omega)=d_{k_n(\omega)}(\omega);
\]
set $z_n=0$ on the complementary event.  The least-index events form a
countable measurable partition.  Consequently $z_n$ is measurable.  It is
also essentially separably valued: outside one null set, the countable family
$(d_k)$ takes values in one separable closed subspace, obtained by closing the
span of their countably many essential separable ranges.  Thus
$z_n\in\cS(X)$.

For $\omega\in\Omega_0$, density in (2), continuity of $T_\omega$, and
$T_\omega u=u$ imply
\[
 \inf_k\rho_k(\omega)=0.
\]
Hence $k_n(\omega)$ exists and
$\rho_{k_n(\omega)}(\omega)<1/n$.  Lemma~\ref{lem:residual} gives
\[
 \norm{z_n(\omega)-u(\omega)}
 \le \frac{C(\omega)}{n}\longrightarrow0,
 \tag{7}
\]
where $C=C_{GC}$ or $C_{HR}$.  Define $u=0$ off $\Omega_0$.
Equation (7) expresses $u$ as an almost-sure norm limit of strong random
variables, so $u\in\cS(X)$, and it is a random fixed point.

It remains to prove uniqueness in the random class.  Let $u,v\in\cS(X)$ be
fixed points.  Under (GC), the cross term equals
$2\norm{u-v}$, but $2c<1$, while both residual terms vanish.  Thus
\[
 \norm{u-v}\le a\norm{u-v},
\]
so $u=v$ almost surely.  Under (HR),
\[
 \norm{u-v}
 \le(\alpha_1+\alpha_4+\alpha_5)\norm{u-v},
\]
and the coefficient is below one.  This proves Theorem~\ref{thm:main}.
\qed

\section{Verification and scope}

\paragraph{Checks performed.}
The proof was audited separately for (i) closure of strong measurability under
composition and countable pasting, (ii) synchronization of null sets using
only the constructed countable family, (iii) invariance of $Y_\omega$, (iv)
positivity of every residual-bound denominator, and (v) uniqueness among
strong random fixed points.  The included script
\texttt{code/verify\_residual\_bounds.py} stress-tested 100,000 admissible
scalar configurations for each of (4) and (5); it returned \texttt{PASS}.
This computation is only a sanity check; Lemma~\ref{lem:residual} is the proof.

\paragraph{Exact scope.}
The theorem removes separability of the Banach space completely.  It does not
assert that the source's bare Borel-section wording is sufficient in every
nonseparable codomain.  In that setting, Borel measurability is not stable
enough for the source's composition steps.  Strong measurability is the
standard nonseparable replacement and is no extra restriction in the
separable setting.  The result therefore answers the source's broad
generalization question affirmatively while making the indispensable
nonseparable measurability convention explicit.

\paragraph{Bounded novelty search.}
On 11 August 2026 we searched the exact source title and question, its citation
neighborhood, and combinations of \emph{Hardy--Rogers}, \emph{nonseparable},
\emph{strongly measurable}, \emph{strong random operator}, \emph{random fixed
point}, and \emph{countable invariant subspace}.  Guo, Zhang, Wang, and Yuan
prove a nonseparable theorem for strong random nonexpansive operators on
weakly compact convex sets with normal structure \cite{guo}; that result
supports the formulation but does not cover the present whole-space
Gregu\v{s}--\'{C}iri\'{c} or Hardy--Rogers theorem.  No exact answer or the
residual-minimizer construction above was located.  This is a bounded search,
not an exhaustive novelty certification.

\paragraph{Human-review recommendation.}
Prioritize the countable-invariant-hull step, the use of the deterministic
Gregu\v{s}--\'{C}iri\'{c} theorem on $Y_\omega$, and the two residual bounds.
Also confirm that the strong-random formulation is accepted as the intended
nonseparable reading of the source question.

\begin{thebibliography}{9}
\bibitem{source}
P.~Saipara, P.~Kumam, and Y.~J.~Cho,
\emph{Random fixed point theorems for Hardy--Rogers self-random operators
with applications to random integral equations}, Stochastics 90 (2018),
297--311; arXiv:1706.01634.

\bibitem{saha}
M.~Saha and A.~Ganguly,
\emph{Random fixed point theorem on a \v{C}iri\'{c}-type contractive mapping
and its consequence}, Fixed Point Theory Appl. 2012, Article 209 (2012).

\bibitem{ciric2000}
L.~\'{C}iri\'{c},
\emph{On a generalization of a Gregu\v{s} fixed point theorem},
Czechoslovak Math. J. 50 (2000), 449--458.

\bibitem{ciric1974}
L.~B.~\'{C}iri\'{c},
\emph{A generalization of Banach's contraction principle},
Proc. Amer. Math. Soc. 45 (1974), 267--273.

\bibitem{guo}
T.~Guo, E.~Zhang, Y.~Wang, and G.~Yuan,
\emph{$L^0$-convex compactness and random normal structure in
$L^0(\mathcal F,B)$}, Acta Math. Sci. Ser. B 40 (2020), 457--469;
arXiv:1904.03607.
\end{thebibliography}

\end{document}
